% Day 5's single deck: what turns quantisation into real speed, before Lab 1.
\begin{frame}{Real Speed Needs Integer Kernels}
    \begin{itemize}
        \setlength\itemsep{0.45em}
        \item \textbf{Fake quantisation} rounds the weights to the integer grid and straight back to floats. The model
              still computes in FP32, so it is no smaller and no faster. It is how you \emph{measure the accuracy
              cost} before deploying anything, and how quantisation-aware training works.
        \item \textbf{Real gains need a runtime with integer kernels:} ONNX Runtime on CPUs, TensorRT on NVIDIA GPUs,
              Core ML on Apple devices, OpenVINO on Intel. You export the model (for example to ONNX), quantise it
              with a calibration set, and run it there.
        \item \textbf{How much faster depends on the chip.} CPUs with INT8 dot-product instructions gain the most.
              Benchmark on the device you will ship to.
        \item \textbf{Not every layer should be quantised.} A tensor that mixes values of very different sizes loses
              the small ones to a single shared scale. Keep such layers in float, and re-measure accuracy after
              every change.
    \end{itemize}
    \vspace{0.3em}
    \begin{block}{In Lab 1}
        You quantise a CNN by hand, then for real with ONNX Runtime, then a YOLO11 object detector, and compare the
        versions live on a webcam.
    \end{block}
\end{frame}
